\documentclass[10pt]{beamer}
\usetheme{metropolis}
\usepackage{graphicx}
\usepackage{listings}
\usepackage{booktabs}
\usepackage{xcolor}
\usepackage{hyperref}

\definecolor{codebg}{HTML}{F5F5F5}
\definecolor{codegreen}{HTML}{2E7D32}
\definecolor{codegray}{HTML}{757575}
\definecolor{codeblue}{HTML}{1565C0}

\lstdefinestyle{rstyle}{
  language=R, backgroundcolor=\color{codebg}, basicstyle=\ttfamily\scriptsize,
  keywordstyle=\color{codeblue}\bfseries, stringstyle=\color{codegreen},
  commentstyle=\color{codegray}\itshape, breaklines=true, frame=single,
  rulecolor=\color{codegray!30}, showstringspaces=false, tabsize=2,
  morekeywords={library,ggplot,aes,geom_col,geom_histogram,geom_point,geom_boxplot,
    geom_line,geom_bar,geom_smooth,geom_text,geom_vline,geom_hline,
    coord_flip,scale_x_log10,scale_color_manual,scale_fill_manual,
    facet_wrap,stat_summary,theme_minimal,theme,element_text,element_blank,
    labs,ggsave,fct_reorder,count,slice_max,filter,mutate,drop_na,
    patchwork,plot_annotation,tag_levels}
}
\lstdefinestyle{pythonstyle}{
  language=Python, backgroundcolor=\color{codebg}, basicstyle=\ttfamily\scriptsize,
  keywordstyle=\color{codeblue}\bfseries, stringstyle=\color{codegreen},
  commentstyle=\color{codegray}\itshape, breaklines=true, frame=single,
  rulecolor=\color{codegray!30}, showstringspaces=false, tabsize=4,
  morekeywords={import,from,as,pd,plt,sns,np,fig,ax,subplots,GridSpec,
    savefig,tight_layout,suptitle,text}
}

\title{Lab: Building a Visual Report}
\subtitle{Workshop 2 --- Module 10}
\author{Prof.Asoc.\ Endri Raco}
\institute{Data Visualization for Data Scientists \\ UNYT Tirana}
\date{2025/26}

\begin{document}

\begin{frame}
  \titlepage
\end{frame}

\begin{frame}{Module Roadmap}
  \begin{enumerate}
    \item The visual report workflow: question to export
    \item Apply all Grammar of Graphics layers to six chart types
    \item Compose a six-panel dashboard in R (patchwork) and Python (GridSpec)
    \item Workshop 2 review: what you now know
  \end{enumerate}
  \begin{exampleblock}{Learning Outcome}
    You will produce a complete, publication-quality, six-panel visual
    report from the Google Play Store dataset, applying every grammar
    layer covered in Modules 1--9.
  \end{exampleblock}
\end{frame}

\section{The Visual Report Workflow}

\begin{frame}{Six Steps: Question to Export}
  \begin{center}
    \includegraphics[width=0.95\textwidth]{figures_w02m10/report_structure}
  \end{center}
  \begin{alertblock}{Pro-Tip}
    Start with the question, not the chart type. ``Which categories
    dominate the Play Store?'' determines the chart (bar); ``How are
    ratings distributed?'' determines the chart (histogram). The grammar
    layers follow from the question.
  \end{alertblock}
\end{frame}

\section{The Target: Six-Panel Dashboard}

\begin{frame}{Complete Visual Report}
  \begin{center}
    \includegraphics[width=0.92\textwidth]{figures_w02m10/dashboard}
  \end{center}
\end{frame}

\begin{frame}{Grammar Layers Applied per Panel}
  \begin{center}
    \includegraphics[width=0.92\textwidth]{figures_w02m10/grammar_checklist}
  \end{center}
\end{frame}

\section{Building Each Panel}

\begin{frame}[fragile]{Panel (a): Top Categories --- Bar Chart}
  \begin{columns}[T]
    \begin{column}{0.52\textwidth}
\begin{lstlisting}[style=rstyle]
# Question: Which categories dominate?
apps_clean |>
  count(Category, sort = TRUE) |>
  slice_max(n, n = 8) |>
  mutate(
    Category = fct_reorder(Category, n),
    highlight = Category == "FAMILY") |>
  ggplot(aes(x = Category, y = n,
             fill = highlight)) +
  geom_col(width = 0.5,
           show.legend = FALSE) +
  geom_text(aes(label = scales::comma(n)),
    hjust = -0.15, size = 3,
    fontface = "bold") +
  scale_fill_manual(values =
    c("TRUE" = "#1565C0",
      "FALSE" = "#BBBBBB")) +
  coord_flip() +
  theme_minimal() +
  theme(panel.grid.major.y =
    element_blank()) +
  labs(title = "FAMILY Leads with 1,832",
       x = NULL, y = NULL)
\end{lstlisting}
    \end{column}
    \begin{column}{0.45\textwidth}
      \includegraphics[width=\textwidth]{figures_w02m10/r_report}

      \vspace{4pt}
      {\small Grammar layers:
      \texttt{geom\_col} + \texttt{coord\_flip}
      + \texttt{scale\_fill\_manual} (accent)
      + \texttt{geom\_text} (direct labels)
      + declarative title.}
    \end{column}
  \end{columns}
\end{frame}

\begin{frame}[fragile]{Panel (a): Same Chart in Python}
  \begin{columns}[T]
    \begin{column}{0.52\textwidth}
\begin{lstlisting}[style=pythonstyle]
top8 = (apps_clean
    .value_counts("Category")
    .head(8)
    .sort_values())

fig, ax = plt.subplots(figsize=(6, 5))
colors = ["#BBBBBB"] * 8
colors[-1] = "#1565C0"

ax.barh(top8.index, top8.values,
    color=colors, height=0.5)

for i, v in enumerate(top8.values):
    ax.text(v + 15, i, f"{v:,}",
        va="center", fontsize=7,
        fontweight="bold",
        color="#1565C0"
            if i == 7 else "#888")

ax.set_title("FAMILY Leads with 1,832",
    fontweight="bold")
for sp in ["top","right","left"]:
    ax.spines[sp].set_visible(False)
ax.set_xticks([])
ax.tick_params(left=False)
\end{lstlisting}
    \end{column}
    \begin{column}{0.45\textwidth}
      \includegraphics[width=\textwidth]{figures_w02m10/py_report}
    \end{column}
  \end{columns}
\end{frame}

\begin{frame}[fragile]{Panels (b)--(f): Grammar Layer Summary}
  \centering
  {\small
  \begin{tabular}{lp{3.5cm}p{3.5cm}}
    \toprule
    \textbf{Panel} & \textbf{R Code Pattern} & \textbf{Python Equivalent} \\
    \midrule
    (b) Histogram & \texttt{geom\_histogram(bins=30)} + \texttt{geom\_vline(mean)} & \texttt{ax.hist(bins=30)} + \texttt{ax.axvline(mean)} \\
    (c) Scatter & \texttt{geom\_point(aes(color=Type))} + \texttt{scale\_x\_log10()} & \texttt{ax.scatter(c=)} + \texttt{set\_xscale("log")} \\
    (d) Boxplot & \texttt{geom\_boxplot(notch=T)} + \texttt{stat\_summary(mean)} & \texttt{ax.boxplot(notch=T)} + \texttt{scatter(means)} \\
    (e) Line & \texttt{geom\_line()} + direct labels & \texttt{ax.plot()} + \texttt{ax.text()} \\
    (f) Stacked & \texttt{geom\_bar(position="stack")} & \texttt{ax.bar(bottom=)} \\
    \bottomrule
  \end{tabular}}
\end{frame}

\section{Composing the Dashboard}

\begin{frame}[fragile]{R: patchwork Composition}
\begin{lstlisting}[style=rstyle]
library(patchwork)

# Compose 6 panels: 3 on top, 3 on bottom
dashboard <- (p_bar | p_hist | p_scatter) /
             (p_box | p_line | p_stacked) +
  plot_annotation(
    title = "Google Play Store: Visual EDA Report",
    subtitle = "Workshop 2 Lab — Grammar of Graphics Applied",
    caption = "Source: Kaggle | UNYT Data Visualization Course",
    tag_levels = "a"  # auto-labels (a), (b), ...
  )

# Export
ggsave("report.pdf", dashboard, width = 14, height = 9, device = cairo_pdf)
ggsave("report.png", dashboard, width = 14, height = 9, dpi = 300)
\end{lstlisting}
\end{frame}

\begin{frame}[fragile]{Python: GridSpec Composition}
\begin{lstlisting}[style=pythonstyle]
from matplotlib.gridspec import GridSpec

fig = plt.figure(figsize=(14, 9))
gs = GridSpec(2, 3, figure=fig, hspace=0.4, wspace=0.35)

ax_bar    = fig.add_subplot(gs[0, 0])   # top-left
ax_hist   = fig.add_subplot(gs[0, 1])   # top-center
ax_scat   = fig.add_subplot(gs[0, 2])   # top-right
ax_box    = fig.add_subplot(gs[1, 0])   # bottom-left
ax_line   = fig.add_subplot(gs[1, 1])   # bottom-center
ax_stack  = fig.add_subplot(gs[1, 2])   # bottom-right

# ... plot each panel onto its axes ...

fig.suptitle("Google Play Store: Visual EDA Report",
    fontsize=13, fontweight="bold")
fig.text(0.5, 0.01, "Source: Kaggle | UNYT",
    ha="center", fontsize=8, fontstyle="italic", color="#888")

fig.savefig("report.pdf", bbox_inches="tight")
fig.savefig("report.png", dpi=300, bbox_inches="tight")
\end{lstlisting}
\end{frame}

\section{Workshop 2 Review}

\begin{frame}{Workshop 2: What You Now Know}
  \begin{center}
    \includegraphics[width=0.85\textwidth]{figures_w02m10/w02_summary}
  \end{center}
\end{frame}

\begin{frame}{The Grammar Gives You Infinite Charts}
  \centering
  \begin{tabular}{rlp{5.5cm}}
    \toprule
    \textbf{Layer} & \textbf{R} & \textbf{Python} \\
    \midrule
    Data + Aes   & \texttt{ggplot(df, aes())}     & \texttt{ggplot(df, aes())} / loop \\
    Geometry     & \texttt{geom\_*()} (20+ types)  & \texttt{ax.*} or \texttt{geom\_*()} \\
    Statistics   & \texttt{stat\_*()} or default   & Built into seaborn / manual \\
    Scales       & \texttt{scale\_*()} family      & \texttt{set\_xscale} / \texttt{cmap} \\
    Facets       & \texttt{facet\_wrap/grid}       & \texttt{FacetGrid} / subplots \\
    Coordinates  & \texttt{coord\_*()} family      & projection / set\_aspect \\
    Theme        & \texttt{theme\_*() + theme()}   & \texttt{rcParams} / manual \\
    Labels       & \texttt{labs()}                 & \texttt{set\_title/xlabel/text} \\
    Annotations  & \texttt{annotate() + geom\_*}   & \texttt{ax.annotate/axhline} \\
    \bottomrule
  \end{tabular}
\end{frame}

\begin{frame}{Key Takeaways}
  \begin{enumerate}
    \item A visual report is built in six steps: \textbf{question → wrangle
          → explore → polish → compose → export}.
    \item Every chart is a stack of \textbf{grammar layers}: data, aes,
          geom, stat, scale, facet, coord, theme, labs, annotations.
    \item \textbf{patchwork} (R) and \textbf{GridSpec} (Python) compose
          individual charts into multi-panel dashboards.
    \item The \textbf{Grammar of Graphics} is a universal system: learn
          it once in ggplot2, apply it in plotnine, seaborn, or matplotlib.
    \item \textbf{Workshop 3} builds on this foundation with the Chart
          Types Encyclopedia: distributions, comparisons, relationships,
          proportions, rankings, and uncertainty.
  \end{enumerate}
\end{frame}

\begin{frame}{Workshop 2 Homework Summary}
  \begin{alertblock}{Due: Before Workshop 3}
    Create a ``Grammar of Graphics Cheat Sheet'': implement 15 chart
    types using GoG layers in both R and Python.
    \begin{itemize}
      \item 15 chart types (bar, histogram, scatter, box, line, area,
            lollipop, dumbbell, heatmap, violin, density, stacked,
            faceted scatter, bubble, annotated line)
      \item Each in R (ggplot2) and Python (plotnine or seaborn)
      \item One-page PDF per chart showing code + output
    \end{itemize}
    Total weight: 5\% of course grade.
  \end{alertblock}
  \vspace{4pt}
  \textbf{Next}: Workshop 3 --- Chart Types Encyclopedia (M01: Distributions)
\end{frame}

\end{document}
